\section{Introduction}

In dynamic algorithms, the main goal is to maintain a key property of the graph while an adversary makes changes to the edges of the graph. An algorithm is called incremental if it only handles edge insertions, and decremental if it handles only edge deletions and fully dynamic if it handles both edge insertions and deletions. \\ \\
The problem of maintaining a $(1-\eps)$-approximation to the maximum matching in a dynamic graph is a well-studied one. In the fully dynamic setting, the best known update time is $O(m^{0.5-\Omega_{\eps}(1)})$ (see \cite{BKS23}), and the conditional lower bounds proved in the \cite{HKNS15,KPP16} suggest that this is a hard barrier to break through. Consequently, several relaxations of this problem have been studied. For example, one line of research has shown that we can get considerably faster update times if we settle for large approximation factors. Another research direction has been to consider more relaxed incremental or decremental models. In the incremental setting, there have been a series of upper and lower bound results (see \cite{Dahlgaard16} and \cite{Gupta2014}), culminating in the results of \cite{GLSSS19} (who give a $O(m\cdot (\nicefrac{1}{\eps})^{O(\nicefrac{1}{\eps})})$ total update time algorithm for $(1-\eps)$-approximate matching in \textbf{general graphs}) and of \cite{BlikstadK23} (who give a $O(m\cdot \poly(1/\eps))$ total update time algorithm for $(1-\eps)$-approximate matching in \textbf{bipartite graphs}). In contrast, for decremental matching \cite{BPT20,JambulapatiJST22} gave an $O(m\cdot\poly(\log n,\nicefrac{1}{\eps}))$ total update time algorithm for $(1-\eps)$-approximate matching in bipartite graphs, and \cite{AssadiBD22} extended this algorithm to the case of general graphs by giving an algorithm for $(1-\eps)$-approximate decremental matching with total time $O_{\eps}(m\cdot \poly(\log n))$, thus completing the picture for approximate cardinality matching in the partially dynamic setting. \\ \\
For the weighted case (in both fully and partially dynamic settings) the picture is less clear. \cite{BernsteinDL21} showed that for the bipartite case, it is possible to reduce weighted matching to unweighted matching with an overhead of $(1-\eps)$ in the approximation factor, and at a loss of a factor of $\log W$ and $(\nicefrac{1}{\eps})^{O(\nicefrac{1}{\eps})}$ in the update time. Recently, \cite{BhattacharyaKS23b} showed that this exponential-in-$1/\eps$ dependence was avoidable for bipartite graphs in the partially dynamic setting. \\ \\
In contrast, for non-bipartite graphs, we don’t have such a general reduction. Thus, while there has been significant recent progress on unweighted matching, the state-of-art for weighted matching lags far behind. Our main contribution in this paper, is to close this gap between weighted and unweighted matching for general graphs in the decremental dynamic setting:

%\begin{theorem2}
%Let $G$ be a weighted graph with weight function $w$ on the edges of the grpah, and let $\eps\in (0,\nicefrac{1}{2})$. Then, there is a deterministic incremental algorithm with total update time $O_{\eps}(m)$ (amortized $O_{\eps}(1)$) that maintains an integral matching $M$ with $w(M)\geqslant (1-\varepsilon)\cdot \text{mwm}(G)$, where $G$ refers to the current version of the graph. 
%\end{theorem2}

%This should be compared to the incremental algorithm for unweighted $(1-\eps)$-approximate matching, which is also deterministic and takes amortized time $\nicefrac{1}{\eps}^{O(\nicefrac{1}{\eps})}$. 

\begin{restatable}{theorem2}{mainthm}\label{thm:main}
Let $G$ be a weighted graph with weight function $w$ on the edges of the graph and let $\eps\in (0,\nicefrac{1}{2})$. Then, there is a decremental algorithm with total update time $O_{\varepsilon}(m\cdot\poly(\log n)\cdot\log R)$ (amortized $\Tilde{O}_{\varepsilon}(1)$) that maintains an integral matching $M$ with $w(M)\geqslant (1-\varepsilon)\cdot \textsf{mwm}(G)$, with high probability, where $G$ refers to the current version of the graph. The algorithm is randomized but works against an adaptive adversary. The dependence on $\eps$ is $2^{{\paren{\nicefrac{1}{\varepsilon}}}^{O\paren{\nicefrac{1}{\varepsilon}}}}$.
\end{restatable}

In comparison, the decremental algorithm for unweighted $(1-\eps)$ approximate matching, is also randomized but works against an adaptive adversary, and takes the same total time albeit without an additional loss of a $\log R$ factor and a dependence of $(\nicefrac{1}{\eps})^{O(\nicefrac{1}{\eps})}$ on $\eps$. 


\subsection{Our Techniques}

Several unweighted dynamic algorithms maintain a matching sparsifier, which is a sparse subgraph of $G$ guaranteed to contain a large matching. As mentioned before, unlike in the bipartite case, there is no reduction from weighted to unweighted matchings in general graphs. So, the existing algorithms for weighted matchings in dynamic setting proceed by looking at weighted analogues of unweighted matching sparsifiers. For example, \cite{BernsteinDL21} consider weighted versions of the EDCS and kernels which were first discovered in the context of cardinality matching problem (see \cite{BS15,BS16,BHI15}). We use a similar approach as in \cite{BernsteinDL21} and port the structural results of \cite{AssadiBD22} for unweighted graphs to the setting of arbitrary weighted graphs. 

%\subsubsection{Structural Result for Incremental Setting}

%The incremental algorithm of [GLSSS19] proceeds by growing alternating path trees to detect new augmenting paths as edges are added to the graph. However, if we want $O_{\eps}(1)$ update time, these trees can’t be too large. Thus, they impose this condition that the trees have short height and bounded degree. Additionally, they also impose the condition that every vertex can appear at a particular height in at most one alternating path tree. The height condition is taken care of because of the following folklore lemma by Hopcraft-Karp:

%\begin{lemma2}(Informally)
% Let $\card{M}$ be the current matching in the graph, then there exist a vertex disjoint set of short (length at most $\nicefrac{1}{\eps}$) augmenting paths, along which we can augment so that the resulting matching $\card{M’}$ has size at least $(1-\eps)\cdot \mu(G)$. 
% \end{lemma2}

%Finally, the bounded degree condition is taken care of by a careful charging argument: suppose we omit a matching edge due to the bounded degree condition at a particular vertex $v$. Due to the disjointness of the augmenting paths there is at most one such edge, and this edge can be charged to $\text{deg}(v)$ many matched edges in the tree containing $v$. Every edge appears at a particular height in at most one tree, and the height is bounded. Thus, choosing the degree to be a large enough constant ensures that the charge on any matched edge is at most $\eps$, and since any missed matched edge can be charged in this way, the total number of missed edges is at most $\eps\cdot \mu(G)$. 

%\subsubsection{Weighted Analog of the Hopcroft-Karp Lemma}

%In the weighted setting, it is no longer the case that we only have to look at augmenting paths. In particular, taking an alternating cycle or path may also increase the weight of the matching. There are also several other obstacles when implementing the above-mentioned algorithm.

%\begin{enumerate}
%\item It is no longer clear that it is sufficient to look at short alternating paths or cycles. 
%\item If we impose the degree condition on the alternating path trees, then the edge we end up omitting may have a substantially higher weight than the edges we wish to charge. Thus, it is no longer possible to maintain a constant degree. 
%\end{enumerate}

%In the case of general graphs, there are also additional impediments due to the presence of blossoms. Thus, we need a weighted analogue of the Hopcroft-Karp Lemma:

%\begin{lemma2}
%Let $M$ be the current matching in the graph, then there exist a vertex disjoint set of short (length at most $\nicefrac{1}{\eps}$) alternating paths, along which we can augment so that the resulting matching $M’$ has $w(M’)\geqslant  (1-\eps)\textsf{mwm}(G)$. Moreover, the ratio of the highest and lowest weight of any edge on any path is bounded. 
% \end{lemma2}

%Thus, we are able to show that there is a set of vertex disjoint augmenting paths of length at most $\nicefrac{1}{\eps}$, which on augmentation give us a matching of near-optimal weight. Finally, since any ``missed" matched edge is not much heavier than the edges to which it has to be charged, we can now afford to build trees of bounded degrees.

\subsubsection{Structural Results for Decremental Matching}

The congestion balancing framework of \cite{BPT20} maintains a fractional matching that is robust to deletions. In order to ensure that the matching is robust and “spread out”, they impose a capacity function $\kappa$ on the edges of the graph. They repeatedly invoke a subroutine $\text{M-or-E*}()$ that does one of the following (here $\mu(G)$ refers to the size of maximum cardinality matching):

\begin{enumerate}
\item Either output a fractional matching $\vec{x}$, such that $\sum_{e\in E}x(e)\geqslant (1-\eps)\cdot \mu(G)$ or,
\item Output a set $E^*$ of bottleneck edges along which we can increase capacities: \label{prop:congestb}
\begin{enumerate}[label=(\roman*)]
\item $\kappa(E^*)$ is $O(\mu\log n)$\label{prop:congesti}
\item Every large matching in $G$ has $\eps\cdot\mu(G)$ edges in $E^*$. \label{prop:congestii}
\end{enumerate}
\end{enumerate} 
The authors of \cite{BPT20} use this subroutine to get an algorithm for decremental matching as follows: either \textsc{M-or-E*}() outputs a fractional matching $\vec{f}$ such that $\sum_{e\in E}\geqslant (1-\eps)\cdot \mu(G)$, in which case we can use a black-box of \cite{Wajc2020} to round this $\vec{f}$, or if there is no $\vec{f}$ satisfying this property, then it outputs a set of edges $E^*$ along which we can increase capacity. Note that, \ref{prop:congestb}\ref{prop:congesti} ensures that the total capacity increase is small, and \ref{prop:congestb}\ref{prop:congestii} ensures that we only increase capacity along edges necessary to form a large matching. Thus, they are able to show that the capacities remain small on average. \\ \\
Predictably in general graphs, implementing such a subroutine is challenging because of the blossom constraints which impose the following roadblocks:
\begin{enumerate}
\item An arbitrary fractional matching in a general graph can suffer from an integrality gap, which is not the case for bipartite graphs. Consequently, we need to be able to check if a given general graph has a large fractional matching obeying capacities and the blossom constraints. \label{item:issuea}
\item Once we know there is a large fractional matching obeying the above requirements, we need to efficiently find such a fractional matching. In bipartite graphs, this can be done use approximate max flows. However, for general graphs this correspondence between fractional matchings and flows is lost. \label{item:issueb}
\item Finally, in the case of bipartite graphs, if the fractional matching is small, the bottleneck edges $E^*$ along which we have to increase capacity correspond to a min-cut in the graph, however for the case of general graphs, this characterization is not clear.\label{item:issuec}
\end{enumerate}

The main contribution of \cite{AssadiBD22} is to address these challenges through three structural lemmas, we state these informally, and then state our weighted versions of these lemmas.

\begin{lemma2}\label{lem:detect}
Let $G$ be a graph and let $\mu(G)$ denotes the size of the optimal matching in $G$. Suppose $\kappa$ is the capacity function on the edges of the graph, and let $\mu(G,\kappa)$ denote the optimal fractional matching of $G$, obeying capacity function $\kappa$, and all odd set constraints. Sample every edge $e$ with $p(e) \propto \kappa(e)$ to create uncapacitated graph $G_s$. Then, $\mu(G_s)\geqslant \mu(G,\kappa)-\eps\cdot n$.
\end{lemma2}

The above lemma gives us a way to determine whether a given capacitated graph has a large fractional matching obeying blossom constraints by reducing this problem to the problem of finding an integral matching in a general graph. For the latter problem, there are efficient algorithms known. This takes care of Issue \ref{item:issuea}

\begin{lemma2}\label{lem:find}
Let $G$ be a graph with capacity function $\kappa$ on the edges of the graph. Let $G_s$ be the graph obtained by sampling edges independently with probability $p(e)\propto \kappa (e)$. Consider a matching $M$ of $G_s$ such that $\card{M}\geqslant (1-\eps)\cdot \mu(G_s)$. Let $M_H\subseteq M$ be the set of edges of high capacity. Let $V_H=V(M_H)$, Let $V_L=V\setminus V_H$. Compute a fractional matching $\vec{f}$ on the low capacity edges of $G[V_L]$, such that $\vec{f}$ obeys $\kappa$. Then, $\card{M_H}+\sum_{e\in E}f(e)\geqslant \mu(G_s)-\eps\cdot n$.
\end{lemma2}

Thus, by using integral matching algorithms, and approximate maximum flow, the above lemma helps us deal with Issue \ref{item:issueb}.

\begin{lemma2}\label{lem:increase}
The set of bottleneck edges of $E^*$ can be identified by considering the edges of $G$ that are not covered by the duals of the matching program of $G_s$. 
\end{lemma2}

The approximate matching algorithms given in the literature, also give us near feasible dual solution (see \cite{DP09}), thus, the above lemma, in combination with these algorithms gives us a way to deal with Issue \ref{item:issuec}. 


\subsection{Weighted Analogues of Structural Results for Decremental Matching}

Similar to the algorithm \textsc{M-or-E*}() in the unweighted decremental setting, we design an algorithm \textsc{WeightedM-or-E*}() that runs in time $O(m)$, and does one of the following two things.

\begin{enumerate}
\item Either output a fractional matching $\vec{x}$, such that $\sum_{e\in E}w(e)\cdot x(e)\geqslant (1-\eps)\cdot \textsf{mwm}(G)$.
\item Output a set E* of bottleneck edges along which we can increase capacities:
\begin{enumerate}[label=(\roman*)]
\item $\sum_{e\in E^*}\kappa(e)\cdot w(e)=O(\textsf{mwm}(G)\cdot\log n)$
\item Every large matching in $G$ has $\eps\cdot \textsf{mwm}(G)$ weight in $E^*$.
\end{enumerate}
\end{enumerate}

The informal statements of the weighted versions of the unweighted lemmas are almost exactly the same:

\begin{lemma2}
Let $G$ be a graph and let $\textsf{mwm}(G)$ denotes the weight of the maximum weight matching in $G$. Suppose $\kappa$ is the capacity function on the edges of the graph, and let $\textsf{mwm}(G,\kappa)$ denote the maximum fractional matching of $G$, obeying capacity function $\kappa$, and all odd set constraints. Sample every edge $e$ with $p(e) \propto \kappa(e)$ to create uncapacitated graph $G_s$. Then, $\textsf{mwm}(G_s)\geqslant \textsf{mwm}(G,\kappa)-\eps\cdot n$.
\end{lemma2}

The proof of \Cref{lem:detect}, relies on the Tutte-Berge Theorem (\cite{Schrijver2003CombinatorialOP}), and the main impediment to proving the above lemma is that there is no weighted analogue of this lemma. Thus, we have to rely on other structural probabilistic tools to come up with a proof.

\begin{lemma2}
Consider a matching M of $G_s$ such that $w(M)\geqslant (1-\eps)\cdot \textsf{mwm}(G_s)$. Let $M_H$ be the set of edges of high capacity. Let $V_H=V(M_H)$, Let $V_L=V\setminus V_H$. Compute a weighted fractional $\vec{f}$ matching on the low capacity edges of $G[V_L]$, such that $\vec{f}$ obeys the capacity function $\kappa$, and the odd set constraints. Then, $\sum_{e\in E}f(e)\cdot w(e)+w(M_H)\geqslant  \textsf{mwm}(G_s)-\eps\cdot n$.
\end{lemma2}

The proof of the unweighted version of this theorem relies on the extended Hall's theorem, and similar to Tutte Berge theorem, there is no weighted version of this theorem. Thus, we have to come up with a generalization that relies on a primal-dual analysis. \\ \\
Note that in the unweighted version, we can easily find such a fractional matching on low capacity edges, using approximate max-flow. However, weighted fractional matchings obeying capacities correspond to min-cost flow, and for this problem, we only have algorithms with $O(m^{1+o(1)})$ time (see \cite{BPT20,BrandCKLPGSS}) or ones where the runtime depends polynomially on the ratio between maximum and minimum capacities, which in our setting can be $O(n^2)$ (see \cite{GoldbergHKT17}). Thus, we have to come up with a simple, efficient algorithm to do this task.

\begin{lemma2}
There is an efficient algorithm to compute $(1-\eps)$-approximate maximum weight fractional matching in a capacitated bipartite graph. 
\end{lemma2}

Finally, we show how to identify the set of bottleneck edges, $E^*$.

\begin{lemma2}The set of bottleneck edges of E* can be identified by looking at the dual of the matching program of $G_s$. 
\end{lemma2}

